\documentclass[12pt,a4paper]{article}

\usepackage[T1]{fontenc}
\usepackage[brazil]{babel}
\usepackage{array}
\usepackage{geometry}
\geometry{margin=2.5cm}

\begin{document}

\begin{center}
    {\Large\bfseries QUÍMICA VIII -- OCEANOGRAFIA}\\[0.3em]
   
    \vspace{0.6em}
    {\bfseries CRONOGRAMA -- PRIMEIRO SEMESTRE / 2026}
\end{center}

\vspace{1em}

\renewcommand{\arraystretch}{1.3}

\begin{center}
\begin{tabular}{|>{\centering\arraybackslash}p{2.2cm}|>{\centering\arraybackslash}p{11cm}|}
\hline
\textbf{DATA} & \textbf{PRÁTICA} \\
\hline
09/03 & Apresentação \\
\hline
16/03 & Técnicas de pesagem (lenta e adição) \\
\hline
23/03 & Calibração de um material volumétrico \\
\hline
30/03 & Preparo de solução de HCl e Bórax \\
\hline
06/04 & Padronização da solução de HCl e Bórax \\
\hline
13/04 & Preparo e padronização de NaOH \\
\hline
20/04 & Recesso \\
\hline
27/04 & P1 -- prova teórica \\
\hline
04/05 & Determinação da alcalinidade em água do mar \\
\hline
11/05 & Determinação de cloreto em água do mar \\
\hline
18/05 & Determinação de Ca e Mg em águas naturais \\
\hline
25/05 & Determinação de O$_2$ dissolvido em água pelo método de Winkler \\
\hline
01/06 & Determinação do teor de SO$_4^{2-}$ em água do mar por gravimetria de precipitação -- Parte 1 \\
\hline
08/06 & Determinação do teor de SO$_4^{2-}$ em água do mar por gravimetria de precipitação -- Parte 2 \\
\hline
15/06 & Aula de dúvidas \\
\hline
22/06 & P2 -- prova teórica \\
\hline
29/06 & Prova de reposição \\
\hline
06/07 & Prova Final -- teórica \\
\hline
\end{tabular}
\end{center}

\vspace{1em}

\noindent\textbf{Início das aulas:} 09/03/26

\noindent\textbf{Término das aulas:} 11/07/26

\end{document}